% 電磁気学1 電磁気学1-3 問題1（001-tani。当時は 電磁気学1-1-2 問題1・shou01-02-l1-tani） 解いた日 2026-10-18
%% kotae
\toi{1}{◎}{}
\[ (\mathrm{s^{2}\,C^{2}\,kg^{-1}\,m^{-3}})\cdot(\mathrm{kg\,m\,C^{-2}})=\mathrm{s^{2}\,m^{-2}} \]
\toi{2}{○}{単位}
\begin{align*}
\frac{1}{\sqrt{\varepsilon_{0}\mu_{0}}}&=2.9979\cdots\times10^{8}\\
&=2.998\times10^{8}\ \nuke{\mathrm{m/s}}\ \nuke{{}=c}
\end{align*}
\naze{数値は正しいが，単位 \unit{\metre\per\second} が書かれていない。また「光速 $c$ と比べよ」への答え（有効数字4桁で $c=\qty{2.998e8}{\metre\per\second}$ と一致する）が無い。}
\toi{3}{◎}{}
\begin{align*}
\frac{1}{4\pi\varepsilon_{0}}&=8.987\cdots\times10^{9}\\
&=8.99\times10^{9}
\end{align*}
\begin{align*}
(\mathrm{s^{2}\,C^{2}\,kg^{-1}\,m^{-3}})^{-1}&=\mathrm{s^{-2}\,C^{-2}\,kg\,m^{3}}\\
&=\mathrm{m^{2}\,C^{-2}}\,(\mathrm{kg\,m\,s^{-2}})\\
&=\mathrm{N\,m^{2}\,C^{-2}}
\end{align*}
%% kizuki
\item \mae{第1-2章 問題1・10-11} 前回は (3) で \unit{\newton} を取り出した残りの指数を落とし，今回は (2) で数値に単位を付け忘れた。どちらも単位の扱いの誤りである。数値の答えを書いたら必ず単位を付け，その単位が速さ（\unit{\metre\per\second}）になっているかを (1) の結果の平方根の逆数で確かめる。
\item 「$c$ と比べよ」のように比較を求められたら，比べた結果（有効数字4桁で $c=\qty{2.998e8}{\metre\per\second}$ と一致する）を一文で書く。計算した値だけでは問いに答えたことにならない。
